% Laget automatisk av elevark/lag_elevark.py fra opplegg/generator-transformator.html. Endringer her blir overskrevet.
\documentclass[11pt]{article}
\usepackage{elevark}
\begin{document}
\elevarkhode{Fysikk 2 · Elektromagnetisk induksjon}{Generatoren, transformatoren og strømnettet – arbeidsark}{https://fysikklab.no/sim/generator.html}{Generator}

\section*{Oppgaver}

Oppgave 1–3 bruker Generator\footnote{\url{fysikklab.no/sim/generator.html}}. Oppgave 4–7 er uten simulering.

\begin{oppgaveliste}
\oppgave \fase{gjett}{Gjett}I hvilken stilling er spenningen størst: når spoleplanet står vinkelrett på feltet, eller når det er parallelt med feltet? Kjør simuleringen og sjekk.
\svar
\oppgave \fase{test}{Regn}Generatoren har $N = 100$, $B = 0{,}20\,\mathrm{T}$, $A = 100\,\mathrm{cm^{2}}$ og $f = 2{,}0\,\mathrm{Hz}$. Regn ut toppspenningen og effektivverdien. Hva skjer når du dobler frekvensen?
\svar
\oppgave \fase{forklar}{Forklar}Kryss av for kommutator. Hva gjør den med spenningen? Hvorfor trenger ikke et kraftverk kommutator?
\svar
\oppgave \fase{test}{Regn}Hvor stor må NBA være for at en generator som går med 50~Hz, skal gi 230~V effektivt?
\svar
\oppgave \fase{test}{Regn}En lader omformer 230~V til 12~V og gir 36~W. Primærspolen har 1150 vindinger. Finn antall vindinger på sekundærspolen og strømmen på hver side (ideell transformator).
\svar
\oppgave \fase{test}{Regn}Et kraftverk sender 100~MW gjennom en linje med resistans 5,0~Ω. Regn ut effekttapet ved 400~kV og ved 20~kV.
\svar
\oppgave \fase{utfordring}{Utfordring}En sykkeldynamo har en spole med 10 vindinger og sidekanter 2,5~cm i et felt på 0,50~T. Den drives av et hjul med diameter 1,5~cm som ruller mot et sykkelhjul med diameter 71~cm. Hvor stor toppspenning gir den når du sykler i 30~km/h?
\svar
\end{oppgaveliste}

\end{document}
